\chapter{Symmetry, Noether's Theorem, and Energy}

Our destination is energy conservation, but first there is a correction to
make. In the previous discussion of motion in polar coordinates, the outward
term in the radial equation was called Coriolis. The correct name is
\emph{centrifugal}. Rather than repair only the word, we shall reconstruct the
whole example. It provides the simplest route from the Euler--Lagrange
equations to symmetry, conservation laws, and finally the Hamiltonian.

\section{A Particle in Polar Coordinates}

The working procedure is straightforward. Choose coordinates that describe
the system, write the kinetic and potential energies, form the Lagrangian, and
apply the Euler--Lagrange equations. The difficult part is usually choosing a
good description; after that, the calculation is largely systematic.

For a particle moving in a plane, Cartesian coordinates \(x,y\) are possible,
but polar coordinates \(r,\theta\) are better adapted to a force directed
toward or away from an origin. The radial velocity is
\begin{equation}
 v_r=\dot r,
\end{equation}
whereas the transverse velocity is
\begin{equation}
 v_\theta=r\dot\theta.
 \label{eq:l04-transverse-speed}
\end{equation}
The factor of \(r\) matters. At the same angular velocity, a point farther from
the origin sweeps out a larger distance in the same time.

The two velocity components are perpendicular, so
\begin{equation}
 v^2=\dot r^{\,2}+r^2\dot\theta^{\,2},
\end{equation}
and the kinetic energy is
\begin{equation}
 T=\frac{m}{2}\left(\dot r^{\,2}+r^2\dot\theta^{\,2}\right).
\end{equation}
We now choose a central potential,
\begin{equation}
 U=U(r),
\end{equation}
and therefore
\begin{equation}
 \mathcal L
 =\frac{m}{2}\left(\dot r^{\,2}+r^2\dot\theta^{\,2}\right)-U(r).
 \label{eq:l04-polar-lagrangian}
\end{equation}
This restriction is not compulsory. A potential may depend on \(\theta\), but
then the system distinguishes one direction from another. Independence from
\(\theta\) is precisely rotational symmetry.

\begin{classroomqa}
\audiencequestion{Does \(U(r)\) also depend on \(\theta\)?}
\lecturerresponse{Not in the central-force problem being studied. Gravity due
to a fixed central body is the standard example: the potential depends on the
distance from the body, not on the direction from which it is approached. A
more general potential \(U(r,\theta)\) is possible, but it would not possess
this rotational symmetry.}
\end{classroomqa}

\section{The Radial Equation}

For any generalized coordinate \(q_i\), define its canonical momentum by
\begin{equation}
 \pi_i=\frac{\partial\mathcal L}{\partial\dot q_i}.
 \label{eq:l04-canonical-momentum}
\end{equation}
The Euler--Lagrange equation may then be written
\begin{equation}
 \frac{d\pi_i}{dt}=\frac{\partial\mathcal L}{\partial q_i}.
 \label{eq:l04-euler-lagrange}
\end{equation}

For the radial coordinate,
\begin{equation}
 \pi_r=\frac{\partial\mathcal L}{\partial\dot r}=m\dot r,
\end{equation}
and hence \(d\pi_r/dt=m\ddot r\). There are two contributions to
\(\partial\mathcal L/\partial r\): one from the potential and one from the
factor \(r^2\) in the angular kinetic energy. Thus
\begin{align}
 \frac{\partial\mathcal L}{\partial r}
 &=\frac{\partial}{\partial r}
   \left(\frac{m}{2}r^2\dot\theta^{\,2}-U(r)\right) \\
 &=mr\dot\theta^{\,2}-\frac{dU}{dr}.
\end{align}
The radial equation is therefore
\begin{equation}
 m\ddot r=mr\dot\theta^{\,2}-\frac{dU}{dr}.
 \label{eq:l04-radial-equation}
\end{equation}

\begin{figure}[H]
\centering
\includegraphics[width=0.84\textwidth]{lecture_04_figure_01.png}
\caption{The completed radial equation on the board at 00:09:55. The term
\(mr\dot\theta^2\) comes from differentiating the angular part of the kinetic
energy, while \(-dU/dr\) comes from the potential.}
\lectureframe{00:09:55}
\label{fig:l04-radial-equation}
\end{figure}

The first term on the right is positive in the outward radial direction. It
is the centrifugal contribution. Nothing was added by hand: it emerged from
the \(r\)-dependence of the kinetic energy. This is why the Lagrangian method
is so efficient. Equations of motion in curvilinear coordinates may look
complicated, but they follow from the same mechanical steps as the Cartesian
equations.

\begin{classroomqa}
\audiencequestion{What does it mean to call the centrifugal contribution
\emph{not a real force}?}
\lecturerresponse{The word \emph{real} adds nothing and should be withdrawn.
The useful statement is only that \(mr\dot\theta^2\) is not obtained from the
derivative of the potential \(U\). It arises from the coordinate dependence
of the kinetic term. Nevertheless it is present on the right-hand side of the
radial equation and is part of the generalized radial force. A change of
coordinates changes the appearance of the equation, not the motion.}
\end{classroomqa}

\section{Rotational Symmetry and Angular Momentum}

The equation for \(\theta\) is even more revealing. Its canonical momentum is
\begin{equation}
 \pi_\theta
 =\frac{\partial\mathcal L}{\partial\dot\theta}
 =mr^2\dot\theta.
 \label{eq:l04-angular-momentum}
\end{equation}
Since \(\theta\) does not occur in Eq.~\eqref{eq:l04-polar-lagrangian},
\begin{equation}
 \frac{\partial\mathcal L}{\partial\theta}=0,
\end{equation}
and the angular Euler--Lagrange equation becomes
\begin{equation}
 \frac{d}{dt}\left(mr^2\dot\theta\right)=0.
\end{equation}
Thus
\begin{equation}
 mr^2\dot\theta=L,
 \qquad L=\text{constant}.
 \label{eq:l04-L-constant}
\end{equation}
The momentum conjugate to the angle is the angular momentum. Its value is not
fixed by the law itself; it is fixed by the initial position and velocity. By
contrast, \(\pi_r=m\dot r\) generally changes because the radial force is not
zero.

\begin{figure}[H]
\centering
\includegraphics[width=0.88\textwidth]{lecture_04_figure_02.png}
\caption{The polar equations assembled at 00:18:05. Rotational symmetry makes
\(mr^2\dot\theta\) constant, and the same board retains the radial equation
that will receive this constant as an input.}
\lectureframe{00:18:05}
\label{fig:l04-angular-conservation}
\end{figure}

Solving Eq.~\eqref{eq:l04-L-constant} for the angular velocity gives
\begin{equation}
 \dot\theta=\frac{L}{mr^2}.
\end{equation}
Substitution into Eq.~\eqref{eq:l04-radial-equation} yields
\begin{align}
 m\ddot r
 &=-\frac{dU}{dr}
   +mr\left(\frac{L}{mr^2}\right)^2 \\
 &=-\frac{dU}{dr}+\frac{L^2}{mr^3}.
 \label{eq:l04-effective-radial-equation}
\end{align}

\begin{figure}[H]
\centering
\includegraphics[width=0.78\textwidth]{lecture_04_figure_03.png}
\caption{The effective radial equation at 00:21:10. Once the angular momentum
is fixed, angular motion contributes the outward term \(L^2/(mr^3)\).}
\lectureframe{00:21:10}
\label{fig:l04-centrifugal-barrier}
\end{figure}

The whole influence of the angular coordinate has now become an outward force
proportional to \(1/r^3\). Far from the origin it is weak; close to the origin
it becomes very strong. It may also be written as the derivative of an
effective potential,
\begin{equation}
 m\ddot r=-\frac{dU_{\mathrm{eff}}}{dr},
 \qquad
 U_{\mathrm{eff}}(r)=U(r)+\frac{L^2}{2mr^2}.
 \label{eq:l04-effective-potential}
\end{equation}
The second term is the centrifugal barrier.

For gravity or the Coulomb force, the attractive force varies as \(1/r^2\),
whereas the centrifugal contribution varies as \(1/r^3\). At sufficiently
small \(r\), nonzero angular momentum therefore wins. This is why an orbiting
planet does not simply crash into the central body. If the motion is exactly
radial, however, then \(L=0\), the barrier disappears, and a collision may
occur.

\begin{classroomqa}
\audiencequestion{Is centrifugal force just inertia trying to carry the
particle along a straight line?}
\lecturerresponse{Yes. In Cartesian coordinates the free particle continues
along a straight line. In polar coordinates that same tendency appears as an
outward term in the radial equation. With nonzero angular momentum a straight
line misses the origin; with zero angular momentum it may pass through it.
The coordinate descriptions differ, but the trajectory is the same.}
\end{classroomqa}

\section{Variations and Stationarity}

We need one piece of notation before turning to the general theorem. Let
\(F\) be a function of many variables \(\alpha_i\). Under small changes
\(\delta\alpha_i\), its first-order change is
\begin{equation}
 \delta F
 =\sum_i\frac{\partial F}{\partial\alpha_i}\,\delta\alpha_i
 +O(\delta\alpha^2).
 \label{eq:l04-first-variation-function}
\end{equation}
If \(\delta F=0\) for every independent choice of \(\delta\alpha_i\), then
\begin{equation}
 \frac{\partial F}{\partial\alpha_i}=0
 \qquad\text{for every }i.
\end{equation}
This is the meaning of stationarity to first order. A minimum is stationary,
but stationarity is the precise property needed in the action principle.

For a trajectory the same shorthand becomes
\begin{equation}
 A[q]=\int_{t_1}^{t_2}\mathcal L(q,\dot q,t)\,dt,
 \qquad \delta A=0.
\end{equation}
The variation of the action is
\begin{equation}
 \delta A
 =\int_{t_1}^{t_2}
 \sum_i\left(
 \frac{\partial\mathcal L}{\partial q_i}\,\delta q_i
 +\frac{\partial\mathcal L}{\partial\dot q_i}\,\delta\dot q_i
 \right)dt.
 \label{eq:l04-first-variation-action}
\end{equation}
This notation is not an additional physical principle. It is a compact way of
collecting all first-order changes.

\section{Continuous Symmetries}

A symmetry is a transformation that leaves the action unchanged. Translating
an isolated system is a symmetry when the action depends only on relative
positions. Rotating a rotationally invariant system is another. In each case
the transformation acts on the entire history, not merely on one particle at
one instant.

We shall study continuous symmetries through infinitesimal transformations,
\begin{equation}
 q_i\longrightarrow q_i+\epsilon f_i(q),
 \qquad
 \delta q_i=\epsilon f_i(q),
 \label{eq:l04-infinitesimal-symmetry}
\end{equation}
where \(\epsilon\) is small. The displacement may depend on the point in
configuration space. A finite translation or rotation can be built from many
small ones, so first-order information is sufficient.

For an infinitesimal counterclockwise rotation in the \(x\)-\(y\) plane,
\begin{equation}
 \delta x=-\epsilon y,
 \qquad
 \delta y=\epsilon x,
 \label{eq:l04-infinitesimal-rotation}
\end{equation}
and therefore
\begin{equation}
 f_x=-y,
 \qquad f_y=x.
\end{equation}
The signs can be checked on the axes. At a point on the positive \(x\)-axis,
\(y=0\), so \(x\) does not change to first order while \(y\) increases. At a
point on the positive \(y\)-axis, \(x\) decreases while \(y\) does not change
to first order.

\begin{classroomqa}
\audiencequestion{Must the small rotation angle be measured in radians?}
\lecturerresponse{Radians make Eq.~\eqref{eq:l04-infinitesimal-rotation} take
this simple form. Another angular unit merely rescales \(\epsilon\) and hence
rescales the generator by a constant; the conservation law is unchanged.}
\end{classroomqa}

\begin{classroomqa}
\audiencequestion{Does a cube have rotational symmetry?}
\lecturerresponse{It has a discrete rotational symmetry: certain rotations,
such as ninety degrees about suitable axes, return it to itself. It is not
invariant under an arbitrarily small rotation. A sphere is. The argument below
concerns continuous symmetries that can be built from infinitesimal
transformations, not isolated discrete transformations.}
\end{classroomqa}

\section{Noether's Theorem}

Begin with a trajectory \(q_i(t)\) that satisfies the equations of motion.
An ordinary least-action variation keeps its endpoints fixed. A symmetry
variation is different: it transforms every point of the trajectory, including
the endpoints. The action still does not change, but now it vanishes for a
different reason. It vanishes because the transformation is a symmetry, not
because it belongs to the fixed-endpoint family used to derive the equations
of motion.

Integrating the second term of Eq.~\eqref{eq:l04-first-variation-action} by
parts gives
\begin{align}
 \delta A
 &={}
 \int_{t_1}^{t_2}\sum_i
 \left[
 \frac{\partial\mathcal L}{\partial q_i}
 -\frac{d}{dt}\left(
   \frac{\partial\mathcal L}{\partial\dot q_i}
  \right)
 \right]\delta q_i\,dt \\
 &\quad+
 \left[
 \sum_i\frac{\partial\mathcal L}{\partial\dot q_i}\,
 \delta q_i
 \right]_{t_1}^{t_2}.
 \label{eq:l04-noether-variation}
\end{align}
The bulk integral vanishes because the original path obeys the
Euler--Lagrange equations. The boundary term does not vanish by an endpoint
condition, because the symmetry moved those endpoints. Instead symmetry tells
us that the full \(\delta A\) vanishes. Consequently,
\begin{equation}
 \left[
 \sum_i\pi_i\,\delta q_i
 \right]_{t_1}^{t_2}=0.
\end{equation}
Substitute Eq.~\eqref{eq:l04-infinitesimal-symmetry} and remove the common
constant \(\epsilon\):
\begin{equation}
 \left[Q\right]_{t_1}^{t_2}=0,
 \qquad
 Q=\sum_i\pi_i f_i(q).
 \label{eq:l04-noether-charge}
\end{equation}
Since the endpoints may be any two times,
\begin{equation}
 \frac{dQ}{dt}=0.
\end{equation}

\begin{figure}[H]
\centering
\includegraphics[width=0.88\textwidth]{lecture_04_figure_04.png}
\caption{The endpoint form of the variational argument at 00:47:50. The bulk
Euler--Lagrange term is crossed out because the original path obeys the
equations of motion, leaving the boundary contribution
whose equality at two times is the conservation law.}
\lectureframe{00:47:50}
\label{fig:l04-noether-board}
\end{figure}

This is Noether's theorem in the form needed here. The quantity \(Q\) is
called the Noether charge. The word \emph{charge} is used broadly: it means a
conserved quantity, not necessarily electric charge. Momentum, angular
momentum, and energy are all Noether charges. The theorem also uses two facts
together: the equations arise from an action principle, and the action has a
continuous symmetry. Invariance of an arbitrary equation, without an action
principle, is not by itself this theorem.

\begin{classroomqa}
\audiencequestion{Is the rotation in this argument a parallel transport?}
\lecturerresponse{No. Here a rotation is simply a transformation of the
coordinates and of the entire trajectory. Parallel transport is a different
geometric construction and is not needed for this derivation.}
\end{classroomqa}

\section{Translation and Rotation Charges}

Consider a closed system of many particles. Under a common translation in the
\(x\)-direction,
\begin{equation}
 \delta x_a=\epsilon,
 \qquad
 \delta y_a=\delta z_a=0
\end{equation}
for every particle label \(a\). Thus \(f_{x_a}=1\), and
Eq.~\eqref{eq:l04-noether-charge} gives
\begin{equation}
 Q_x=\sum_a\pi_{x_a}=\sum_a m_a\dot x_a
 =\sum_a p_{x_a}.
 \label{eq:l04-total-momentum}
\end{equation}
The conserved Noether charge is total momentum in the translation direction.
The detailed internal forces were never needed. We required only that moving
the whole isolated system leave its potential energy unchanged. A fixed
external body would spoil that symmetry unless it too were included in the
dynamical system.

\begin{figure}[H]
\centering
\includegraphics[width=0.76\textwidth]{lecture_04_figure_05.png}
\caption{Translation along \(x\) at 00:57:30. Every particle has
\(\delta x_a=\epsilon\), while the conserved charge is the summed
\(x\)-momentum \(\sum_a m_a\dot x_a\).}
\lectureframe{00:57:30}
\label{fig:l04-translation-charge}
\end{figure}

For a rotation in the \(x\)-\(y\) plane, use
Eq.~\eqref{eq:l04-infinitesimal-rotation}. For one particle,
\begin{align}
 Q_{\mathrm{rot}}
 &=p_x(-y)+p_yx \\
 &=xp_y-yp_x.
 \label{eq:l04-angular-charge}
\end{align}
This is the component of \(\mathbf r\times\mathbf p\) perpendicular to the
plane. For many particles,
\begin{equation}
 L_z=\sum_a\left(x_ap_{y_a}-y_ap_{x_a}\right).
\end{equation}
Rotations about the other axes give the other components of the angular
momentum vector.

\begin{figure}[H]
\centering
\includegraphics[width=0.76\textwidth]{lecture_04_figure_06.png}
\caption{The rotational Noether charge at 01:04:10. The board identifies
\(\sum_a(x_ap_{y_a}-y_ap_{x_a})\) with angular momentum.}
\lectureframe{01:04:10}
\label{fig:l04-rotation-charge}
\end{figure}

Once a candidate symmetry is known, constructing its charge is mechanical:
find the infinitesimal generator \(f_i\), multiply by the corresponding
canonical momenta, and sum. Finding all the symmetries of a complicated system
may be difficult; using Noether's formula after finding one is not.

\section{Time Translation}

Suppose the laws and external conditions have no explicit dependence on time.
If one solution is begun a little later, the displaced history is again a
solution. An orbit does not care whether we label a particular state as
eleven o'clock or one minute later; the whole orbit is merely shifted in
time. This is time-translation invariance.

The qualification about external conditions is essential. If we study only
part of a system while a massive external body moves according to a prescribed
schedule, the forces on our subsystem depend explicitly on time. Translating
only the subsystem's history is then not a symmetry.

Let the original trajectory be \(q_i(t)\). Pushing it forward in time by
\(\epsilon\) gives the new function
\begin{equation}
 q_i^{(\epsilon)}(t)=q_i(t-\epsilon).
 \label{eq:l04-time-shifted-path}
\end{equation}
At a fixed value of \(t\), the induced displacement in configuration space is
\begin{align}
 \delta q_i(t)
 &=q_i(t-\epsilon)-q_i(t) \\
 &=-\epsilon\dot q_i(t)+O(\epsilon^2).
 \label{eq:l04-time-variation}
\end{align}

\begin{figure}[H]
\centering
\includegraphics[width=0.82\textwidth]{lecture_04_figure_07.png}
\caption{The time-shifted trajectory at 01:23:30. The board keeps both views
of the transformation together: \(q(t)\mapsto q(t-\epsilon)\) and the induced
fixed-time displacement \(\delta q=-\epsilon\dot q\).}
\lectureframe{01:23:30}
\label{fig:l04-time-shift-geometry}
\end{figure}

\begin{classroomqa}
\audiencequestion{Why does \(\delta q_i=-\epsilon\dot q_i\) contain a minus
sign?}
\lecturerresponse{The whole curve is moved forward in time, but at a fixed
time \(t\) the new curve samples the old curve at the earlier time
\(t-\epsilon\). If \(q\) was increasing, moving the graph upward shifts its
intersection with a fixed horizontal time slice toward smaller \(q\). Drawing
the trajectory with time vertical makes the sign immediate. If the slope is
reversed, then \(\dot q<0\) and the same formula reverses the displacement.}
\end{classroomqa}

\section{The Endpoint Overhangs}

Time translation differs from an ordinary configuration-space symmetry
because it also shifts the interval over which the action is evaluated. Draw
the original and translated trajectories with time vertical. Over most of
the interval they overlap in time, but one has a short extra segment at the
upper end and the other has one at the lower end. These are the overhangs.

On the common interval, the same integration by parts as before gives the
endpoint contribution
\begin{equation}
 \left[
 \sum_i\pi_i\delta q_i
 \right]_{t_1}^{t_2}
 =-\epsilon
 \left[
 \sum_i\pi_i\dot q_i
 \right]_{t_1}^{t_2}.
 \label{eq:l04-time-boundary-term}
\end{equation}
The bulk term again vanishes on the equations of motion.

The upper overhang lasts for the small time \(\epsilon\), so its action is,
to first order,
\begin{equation}
 A_{\mathrm{upper}}=\epsilon\mathcal L(t_2)+O(\epsilon^2).
\end{equation}
Similarly, the lower overhang contributes
\begin{equation}
 A_{\mathrm{lower}}=\epsilon\mathcal L(t_1)+O(\epsilon^2).
\end{equation}
This approximation assumes that \(\mathcal L\) is regular and changes little
over a sufficiently short interval. No singular behavior is being allowed at
the endpoints.

Time-translation symmetry says that the total action difference is zero.
Combining the ordinary boundary term with the two overhangs gives
\begin{equation}
 0=\epsilon
 \left[
 \mathcal L-\sum_i\pi_i\dot q_i
 \right]_{t_1}^{t_2}.
 \label{eq:l04-time-noether-endpoints}
\end{equation}
Therefore
\begin{equation}
 \mathcal L-\sum_i\pi_i\dot q_i
 =\text{constant}.
\end{equation}

\begin{classroomqa}
\audiencequestion{Why did an extra \(\mathcal L\) appear here but not in the
spatial translation or rotation examples?}
\lecturerresponse{A spatial transformation moved the path in configuration
space while the times \(t_1,t_2\) remained fixed. A time translation moves
those temporal endpoints as well. The short overhangs contribute
\(\mathcal L\epsilon\) at the two ends, and their difference supplies the new
term.}
\end{classroomqa}

\section{The Hamiltonian Is the Energy}

The conserved expression has first appeared with the opposite sign from the
usual convention. The negative of a conserved quantity is also conserved, so
we define
\begin{equation}
 H=\sum_i\pi_i\dot q_i-\mathcal L.
 \label{eq:l04-hamiltonian-definition}
\end{equation}
This is the Hamiltonian, and it is the Noether charge associated with time
translation.

For a one-dimensional particle with
\begin{equation}
 \mathcal L=\frac12m\dot x^2-U(x),
\end{equation}
the canonical momentum is \(\pi_x=m\dot x\), and therefore
\begin{align}
 H
 &=\pi_x\dot x-\mathcal L \\
 &=m\dot x^2-\left(\frac12m\dot x^2-U(x)\right) \\
 &=\frac12m\dot x^2+U(x).
 \label{eq:l04-H-T-plus-U}
\end{align}
Thus, for the familiar Lagrangian \(T-U\), the Hamiltonian is the familiar
energy \(T+U\).

\begin{figure}[H]
\centering
\includegraphics[width=0.84\textwidth]{lecture_04_figure_08.png}
\caption{The Hamiltonian definition and its one-dimensional check at
01:33:20. The upper line shows \(H=\sum_i\dot q_i\pi_i-\mathcal L\); the lower
calculation reduces it to kinetic plus potential energy.}
\lectureframe{01:33:20}
\label{fig:l04-hamiltonian-board}
\end{figure}

The deeper definition is not simply \emph{kinetic plus potential}. Energy is
the Noether charge generated by time translation. In ordinary mechanics that
charge reduces to \(T+U\), but Eq.~\eqref{eq:l04-hamiltonian-definition} also
applies when the Lagrangian has a less familiar form.

\begin{classroomqa}
\audiencequestion{Is this circular? Was the Lagrangian not defined using
energy in the first place?}
\lecturerresponse{The general definition of the Lagrangian is the integrand
of the action. For the elementary system we happened to know the convenient
form \(\mathcal L=T-U\), but that is not the universal definition. Once a
Lagrangian is known, Eq.~\eqref{eq:l04-hamiltonian-definition} determines the
Hamiltonian. For the conventional form it reproduces \(T+U\), which is a
check rather than an assumption.}
\end{classroomqa}

\subsection{Regularity and explicit time dependence}

Not every arbitrary function is equally useful as a Lagrangian. At this
stage, two conditions are worth recording. It should depend locally on the
coordinates and velocities, and the momenta
\begin{equation}
 \pi_i=\frac{\partial\mathcal L}{\partial\dot q_i}
\end{equation}
should normally determine the velocities. Locally, the standard regularity
condition is
\begin{equation}
 \det\left(
 \frac{\partial^2\mathcal L}
 {\partial\dot q_i\,\partial\dot q_j}
 \right)\neq0.
 \label{eq:l04-regular-lagrangian}
\end{equation}
This is what permits the passage from Lagrangian to Hamiltonian variables. A
full treatment belongs to Hamiltonian mechanics.

If the Lagrangian depends explicitly on time,
\begin{equation}
 \mathcal L=\mathcal L(q,\dot q,t),
\end{equation}
then shifting a trajectory in time need not leave the action invariant. The
derivation above no longer gives a conservation law, and energy need not be
conserved. The precise evolution law for that case will be developed after
the Hamiltonian formalism is in place.

\begin{classroomqa}
\audiencequestion{How is a Lagrangian obtained from observation rather than
guessed?}
\lecturerresponse{Historically, classical equations of motion were often
found first and only later recognized as Euler--Lagrange equations. The
Lorentz force and Maxwell's equations are examples of laws whose action
formulations were identified after the physical equations had been assembled
from experiment. In modern particle physics, measured scattering processes
are tied more directly to interaction terms in a Lagrangian. In either case
the Lagrangian is constrained by observation and symmetry; it is not pulled
from the air.}
\end{classroomqa}

The compactness of the action is one reason this language became fundamental.
The polar-coordinate equations contain centrifugal terms that would be awkward
to guess directly, while the Lagrangian in
Eq.~\eqref{eq:l04-polar-lagrangian} is easy to remember. The same economy
persists in field theory, electromagnetism, relativity, and quantum theory.
Lagrange's formulation came first; Hamilton's reformulation followed in the
nineteenth century and became the natural bridge to quantum mechanics.

The next practical example will be the harmonic oscillator. At the same time,
the formal development will pass from \((q_i,\dot q_i)\) to
\((q_i,\pi_i)\), where the Hamiltonian becomes the central dynamical object.

\section{Symmetry and Dynamics in One View}

The polar-coordinate correction already contains the pattern of the whole
chapter. Rotational invariance removes \(\theta\) from the Lagrangian, so its
conjugate momentum
\begin{equation}
 L=mr^2\dot\theta
\end{equation}
is conserved. Substituting that constant into the radial equation exposes the
centrifugal barrier \(L^2/(mr^3)\).

The general pattern is Noether's theorem. An infinitesimal symmetry
\(\delta q_i=\epsilon f_i(q)\) has the conserved charge
\begin{equation}
 Q=\sum_i\pi_i f_i(q).
\end{equation}
Translations give momentum and rotations give angular momentum. Time
translation requires one extra piece of bookkeeping because the integration
interval moves. Its endpoint overhangs combine with the ordinary variational
boundary term to give
\begin{equation}
 H=\sum_i\pi_i\dot q_i-\mathcal L.
\end{equation}
The Hamiltonian is therefore not merely a convenient formula. It is the
conserved charge that expresses the independence of the laws of physics from
when an isolated experiment is performed.
